\section{Threats To Validity}\label{sec:threats}

This section assesses the internal, external, construct and conclusion validity of the study's design and results.

\subsection{Internal Validity}
% did x cause y or something else
The comparative baseline in this study is a standalone RocketPy simulation in which the flight logic, including the PI controller, apogee detection and parachute deployment, is directly reimplemented into the controller callback. Any difference between this reimplementation and the F Prime flight software would cause the measured apogee difference to be due to an implementation discrepancy rather than a pure coupling strategy effect. The baseline was constructed to match the F Prime logic as closely as possible, but the two are not formally equivalent.


\subsection{External Validity}
% do results generalise beyond this study
Aside from the informal check against a second rocket configuration in Section \ref{coupling-results}, all experiments were conducted on a single fixed rocket configuration, with a fixed motor, geometry and target apogee of 3000\,m. The quantitative results observed in this study, such as the overshoot in non-blocking strategies, wall-clock overhead, and fault detection timing, are specific to this configuration. Evaluating RITL across rockets with different coasting durations, apogee targets, and aerodynamic profiles would be needed to characterise how the accuracy of the coupling strategy scales across mission types.

% The fixed 10\,Hz callback rate used in RocketPy across all experiments is a design choice, not a result. The one-tick delay identified in Section~\ref{s:evaluation} is a function of this rate, as the magnitude of overshoot will vary at different callback frequencies. The conclusions about the Snapshot and Rate Group accuracy are thus specific to this configuration.

The controller callback rate and the Rate Group tick frequency were both varied experimentally in this study on the same fixed rocket configuration used throughout. The specific error magnitudes reported at each rate are therefore not guaranteed to generalise to different rockets. The sensor model sampling rate and the parachute trigger rate were additionally held fixed throughout and were not varied experimentally, so whether they exhibit a similar relationship remains untested.

The fault injection experiment evaluates a single fault type: a barometric freeze at $t=$10\,s. Other fault modes, such as sensor drift, noise or dropout, could trigger different recovery paths or expose failures not addressed in this study.

The HIL target in this study is a Raspberry Pi 4 running standard RPi OS, without real-time scheduling guarantees. While F Prime supports deployment on bare-metal and RTOS-based microcontrollers, the timing behaviour observed in HIL reflects the behaviour of that platform. Teams deploying on different embedded targets, such as STM32 or similar, may observe different results. The communication layer further limits this generalisability. The Orchestrator and FSW communicate over TCP, routed over Ethernet in HIL, chosen since Lockstep requires reliable delivery. Real flight computers typically read sensors over SPI or UART, unacknowledged communication operating at lower latency than TCP over Ethernet. The wall-clock times and timing variance reported in HIL are therefore specific to this transport configuration, and a lower-latency, unacknowledged link would likely reduce both wall-clock overhead and the practical impact of the one-tick delay.



\subsection{Construct Validity}
% did we measure the right thing

The primary measures of coupling accuracy in this study are apogee error and deployment timings, which are measured relative to the standalone RocketPy baseline. This has two limitations. First, the reference is itself a model, so the apogee differences reported in this study are relative to a modelled trajectory, not a real flight, and the true accuracy of any coupling strategy against an actual flight using the same FSW remains uncharacterised.

Second, while this study adopts a 12\,m acceptance threshold derived from the Camões configuration, acceptable tolerances beyond this bound depend on the specific rocket, target apogee and mission requirements. The results reported in this paper characterise the error introduced by each coupling strategy, but whether that error is significant is mission-dependent.

\subsection{Conclusion Validity}

%do the statistics support the claims
Each configuration was evaluated over 10 runs. For Lockstep and Snapshot strategies, where variance is low, this is sufficient to support the results drawn. For the Rate Group strategy in HIL, which produces a standard deviation of 7.46\,m at 25\,Hz, 10 runs limit the coverage of the distribution. The worst-case apogee error under such timing non-determinism could exceed the observed mean, and a larger sample would be required to characterise this with more confidence. The solver time step sensitivity check is an exception, evaluated with a single run per step size given its exploratory scope. This is discussed further in Section \ref{limitations}.
